\chapter{Introduction to Vectors}
\section{Vectors and Linear Combinations}

\begin{itemize}
  \item \textbf{Linear combination:} Forms like \(3\bm{v}+5\bm{w}\), \(\bm{v}\) and \(\bm{w}\) are vectors.
\end{itemize}

\subsection{Why some equations fail}

For example $\begin{smallmatrix} c+2d=1\\ c+3d=0\\ c+4d=0 \end{smallmatrix}$.
It's actually a linear combination $
  c\begin{bsmallmatrix} 1 \\ 1 \\ 1 \end{bsmallmatrix}
  +d\begin{bsmallmatrix} 2 \\ 3 \\ 4 \end{bsmallmatrix}
  =\begin{bsmallmatrix} 1 \\ 0 \\ 0 \end{bsmallmatrix}
$,
because $\begin{bsmallmatrix} 1 \\ 0 \\ 0 \end{bsmallmatrix}$ is not on that plane.

\section{Lengths and Dot Products}
\subsection{Pythagoras theorem for a 3D Vector}

To find the length of a 3D vector, we apply the Pythagorean theorem twice. For example, consider the vector
\((1,2,3)\). First, calculate the length of its projection onto the \(xy\)-plane, which is \((1,2,0)\). By the
Pythagorean theorem, the length of this base is \(\sqrt{1^2 + 2^2} = \sqrt{5}\).

Next, because this base vector is perpendicular to the vertical component \((0,0,3)\), you can apply the theorem a
second time. The total length of the vector, representing the diagonal of the box, is therefore \(\norm{\bm{v}} =
\sqrt{(\sqrt{5})^2 + 3^2} = \sqrt{5+9} = \sqrt{14}\).

\begin{figure}[H]
  \centering
  \incfig{chapter_01_1.1}
  \caption{The length of a three-dimensional vector.}
\end{figure}

\subsection{The Angle Between Two Vectors}
\phantom{}
\begin{proof}[The dot product is $\bm{v}\cdot\bm{w}=0$ when $\bm{v}$ is perpendicular to $\bm{w}$.]
  When $\bm{v}$ and $\bm{w}$ are perpendicular, they form two sides of a right triangle. The third side is $\pm(\bm{v}-\bm{w})$
\end{proof}

\begin{figure}[H]
  \centering
  \incfig{chapter_01_1.2}
  \caption{Perpendicular vectors}
\end{figure}
